\begin{abstract}
Ask a language model for the move that pays off eight steps from now, and it may give you the move that pays off right away. Existing planning benchmarks cannot tell these apart, because they change the problem whenever they change the horizon. We built a test that holds everything fixed except the stated horizon. Each item is a frozen state from a turn-based card game. Between prompts, only the number of decisions ahead at which the goal is measured changes. Exhaustive search gives the provably best move for every horizon, and the scoring gives no credit for ignoring the instruction. Across eight open models from four families, only the two largest reasoning models change their choices in the right direction. Even they mostly trade the greedy move for a different wrong move. Only about one in three of their changed answers is the best long-horizon move, and their scores do not improve. Models without explicit reasoning ignore the horizon altogether, including a 70-billion-parameter model. An exact rulebook or hidden card names leave this picture unchanged. Thinking longer, in short, is not the same as planning further. The test needs only a deterministic simulator, so it carries over to any domain that has one. We release the states, the exact answers and every model response.
\end{abstract}
